\section{Numerical integration of the Wilson flow}

The discussion in sect.~1 now suggests
to combine the HMC algorithm with
the field transformations generated by the
trivializing flow constructed in the previous section.
In particular, if the Wilson gauge action is used,
the transformations generated by the Wilson flow
may lead to an algorithm with improved sampling efficiency.


\subsection{Forward integration}

There is a wide range of numerical integration methods
that can in principle be used
to integrate the Wilson flow
(see ref.~\cite{HairerEtAl}, for example).
The Euler scheme discussed in the following
performs the integration in time steps of size $\eps$ and
updates the link variables one after another
according to
\begin{equation}
   U(x,\mu)\to U'(x,\mu)=\rme^{\eps [Z(U)](x,\mu)}U(x,\mu),
\end{equation}
where
\begin{equation}
   [Z(U)](x,\mu)=T^a\du^a_{x,\mu}\Wl_0(U).
\end{equation}
Starting from the gauge field $U_t$ at time $t$,
the field at time $t+\eps$
is thus obtained by
running through all links $(x,\mu)$ on the lattice
and updating the link variable residing there.
Note that the ordering of the links matters, since
the old value of $U(x,\mu)$ is replaced by the
new one before going to the next link.

The generator of the flow is explicitly given by
\begin{multline}
   [Z(U)](x,\mu)=
   -\sum_{\nu\neq\mu}\Proj\bigl\{
   U(x,\mu)U(x+\hat{\mu},\nu)U(x+\hat{\nu},\mu)^{-1}U(x,\nu)^{-1}\\
   +U(x,\mu)U(x+\hat{\mu}-\hat{\nu},\nu)^{-1}U(x-\hat{\nu},\mu)^{-1}
   U(x-\hat{\nu},\nu)\bigr\},
\end{multline}
where $\hat{\mu}$ denotes the unit vector in direction $\mu$ and
\begin{equation}
   \Proj\bigl\{M\bigr\}=\frac{1}{2}\bigl(M-M^{\dagger}\bigr)-
   \frac{1}{6}\tr\bigl(M-M^{\dagger}\bigr)
\end{equation}
projects any $3\times3$ matrix $M$ to $\Lie$.
The Euler integration of the Wilson flow thus
amounts to applying a number of ``stout smearing'' steps~\cite{Stout},
except that the link variables are here updated one by one rather than
all at once.


\subsection{Backward integration}

The application of $n$ Euler sweeps
maps the initial field
$V=U_0$ to the field $U=U_{n\eps}$ at time $t=n\eps$. If $t$ is held
fixed and $n$ is taken to infinity, this map converges to the
transformation obtained by integrating the Wilson flow exactly.
However, the HMC algorithm may potentially be combined
with the map defined by the Euler integrator at fixed $n$ and $\eps$.
For this proposition to be a viable option,
the transformation must be invertible, i.e.\ one must
be able to trace back the Euler integration
by inverting the link
update steps one by one in the reverse order.

The question is thus whether
eq.~(5.1) has a unique solution $U(x,\mu)$ given $U'(x,\mu)$
(and keeping all other link variables fixed).
As explained in appendix D,
the answer is affirmative,
for arbitrary values of the field variables, if
\begin{equation}
  |\eps|<\frac{1}{8}.
\end{equation}
Moreover, in this range of $\eps$,
the solution $U(x,\mu)=\rme^{-\eps X_{*}}U'(x,\mu)$
can be obtained through the fixed-point iteration
\begin{align}
  &X_0=0,
  \\[1.0ex]
  &X_{n+1}=\bigl\{[Z(U)](x,\mu)\bigr\}_{
  U(x,\mu)=\rme^{-\eps X_n}U'(x,\mu)},
  \qquad
  n=0,1,2,\ldots,
\end{align}
which converges to $X_{*}$ at an exponential rate.

In appendix D it is also shown that the
Jacobian of the transformation (5.1)
is strictly positive in the range (5.5).
The field transformations obtained through the Euler
integration of the Wilson flow are therefore orientation-preserving
diffeomorphisms of the field manifold, as required
for acceptable maps of field space.


\subsection{Jacobian matrix of the Euler integrator}

The Euler step (5.1) amounts to applying the field transformation
\begin{equation}
  [\euler{x,\mu}(U)](y,\nu)=
  \begin{cases}
    \rme^{\eps [Z(U)](x,\mu)}U(x,\mu)
    & \text{if $(y,\nu)=(x,\mu)$,}\\[1.0ex]
    U(y,\nu) & \text{otherwise,}
  \end{cases}
\end{equation}
to the current gauge field. An Euler sweep is then the composition
product of these transformations over all links $(x,\mu)$.
It is straightforward to show that the Jacobian matrix of
a composed transformation is the product of the Jacobian matrices of
the factors. The Jacobian matrix of the
Euler integrator is therefore
an ordered product of the Jacobian matrices
\begin{equation}
  [\euler{x,\mu,*}(U)](y,\nu;z,\rho)^{ac}
  =-2\,\tr\Bigl\{\du^c_{z,\rho}[\euler{x,\mu}(U)](y,\nu)
  [\euler{x,\mu}(U)](y,\nu)^{-1}T^a\Bigr\}
\end{equation}
of the one-link transformations (5.8).

The matrix (5.9) can be expressed through the derivative
\begin{equation}
  [Z_{*}(U)](y,\nu;z,\rho)^{bc}=
  \du^c_{z,\rho}[Z(U)]^b(y,\nu)
\end{equation}
of the generator of the Wilson flow.
Explicitly, one finds that
\begin{multline}
  [\euler{x,\mu,*}(U)](y,\nu;z,\rho)^{ac}=
  \delta^{ac}\delta_{yz}\delta_{\nu\rho}
  +\delta_{xy}\delta_{\mu\nu}
  \Bigl\{\left(\rme^{\Ad X}-1\right)^{ac}\delta_{yz}\delta_{\nu\rho}\\
  +\eps J(-X)^{ab}[Z_{*}(U)](y,\nu;z,\rho)^{bc}\Bigr\}_{X=\eps[Z(U)](x,\mu)},
\end{multline}
where use was made of the $\Group$ formulae
listed in appendix A and B.


\subsection{Jacobian of the integrated transformations}

The Euler integrator generates a sequence of fields
\begin{equation}
  V=U_0\to U_{\eps}\to U_{2\eps}\to\ldots\to U_{n\eps}=U
\end{equation}
by sweeping through the lattice $n$ times and
updating the link variables one by one in a specified order.
In the following, the intermediate field configurations
obtained starting from $U_{k\eps}$
and updating the link variables
on all links that come before $(x,\mu)$ will be denoted by
$U_{k\eps,[x,\mu]}$.
In particular, $U_{k\eps,[x,\mu]}=U_{k\eps}$ if $(x,\mu)$ is the first link
and
\begin{equation}
  U_{k\eps,[y,\nu]}=\euler{x,\mu}(U_{k\eps,[x,\mu]})
\end{equation}
if $(y,\nu)$ follows $(x,\mu)$ in the chosen link order.

Since the transformation $V\to U=\transne(V)$
is a composition product of one-link update steps,
its Jacobian
\begin{equation}
  \det\transnes(V)=
  \prod_{k=0}^{n-1}
  \prod_{x,\mu}\det\euler{x,\mu,*}(U_{k\eps,[x,\mu]})
\end{equation}
factorizes into the product of the Jacobians of the steps.
The latter coincide with the determinants of certain real
$8\times8$ matrices given explicitly
in appendix C.
